\section{项目总结与展望}
\label{sec:conclusion}

\subsection{项目完成情况}

IronBuddy 最终形成的是一套模块边界清晰、可调试可复测的嵌入式健身教练系统。它把视觉
姿态估计、肌电传感、FSM 计数、GRU 动作质量判断、语音交互、Web 展示和
后台记录接口放在同一套系统里，围绕“动作是否标准、目标肌是否真正参与、
用户如何自然交互”建立了一条较完整的原型主线。

项目的主要收获不只是某个模型或某块硬件，而是多个模块之间的工程整合。
视觉模块解决外部姿态，肌电模块补充内部发力，状态机和 GRU 分别处理
确定性计数与模糊质量判断，语音和 LLM 让系统能以更自然的方式给用户反馈，
端到端调试系统和数据库则让系统维护与长期记录有了入口。

\subsection{仍然存在的不足}

项目也有明显不足。首先，肌电数据规模仍然有限，尤其是跨用户泛化能力不能
只根据本地训练表现下结论。贴片位置、皮肤状态和不同人的肌肉基础值都会
影响输入分布，真实 ESP32 UDP 连续输入也需要继续排查和复测。

其次，语音链路虽然已经做了状态化整理，但最终效果仍依赖现场麦克风距离、
网络状态、ASR 质量和 TTS 播报时机。工具调用、自动疲劳触发和 turn 去重
已经有代码与离线测试支撑，但仍需要完整手测后才能写成稳定结论。

再次，当前硬件形态更偏原型验证。腰包、外接线和贴片方便调试，却不是最终
产品形态。如果继续往产品化方向走，需要考虑贴片固定、小型化电源、外壳
结构、续航、无线稳定性和安全走线。

\subsection{后续方向}

后续可以从三个方向继续推进。

\begin{enumerate}
    \item \term{硬件集成与穿戴优化}：把采集板、ESP32 主控和电源集成到
    一块自制集成板中，通过更短线路提升装配便利性和佩戴安全性。运动服、
    腰包和肩带也可以进一步配合走线，让设备更适合真实训练动作。

    \item \term{扩展动作库和数据集}：在深蹲和弯举之外，继续覆盖硬拉、
    推举、划船等动作。每扩展一个动作，都需要重新确认目标肌和代偿肌、
    贴片位置、动作阶段和训练数据。

    \item \term{完善测试和长期记录}：把手测表、训练脚本、端到端调试系统、
    数据库记录和后台总结结合起来，形成固定的复测流程。
    这样每次改语音、模型或前端，都能判断是否影响主线演示。
\end{enumerate}

\subsection{个人体会}

IronBuddy 是我第一次完整经历从前端到后端、从软件到硬件的端到端项目。
之前的课程作业大多围绕单点问题展开：写一段算法、调一块电路、做一个
Demo 页面就算完成；而这次我必须同时面对完全不同性质的工程对象——
模拟前端的肌电采样和噪声、ESP32 的固件与无线链路、RK3399ProX 板端的
多进程与共享内存、本地 NPU 的 RKNN 量化模型、云端 GPU 的 RTMPose 推理
服务、Flask Web 服务和 PWA 前端、百度 AipSpeech 语音守护进程、DeepSeek
LLM 直连，以及 SQLite 长期记录。任何一个环节都不只是"能不能跑通"
的问题，更涉及"它在出错时长什么样、能不能被定位、改动之后会不会
破坏其他链路"。

真正实践之后才发现，跑通一个嵌入式全栈项目需要进行的 debug 数量是
超出想象的：贴片在运动中会移位、ALSA 在板端掉电后默认静音、子进程
必须用 bracket-trick 才能干净地 pgrep 出来、UDP 偶尔丢包却比 BLE
重传更适合连续传感、状态文件必须 write-then-rename 才能避免读到
半截 JSON、RKNN 量化模型的置信度阈值要从 0.5 降到 0.1 才适配板端
NPU、语音播报和麦克风录入需要状态机隔离才不会自录成回声。这些坑
没有一条来自单门课程，但合起来构成了嵌入式全栈系统能否真正稳定运行
的关键。

把这些散落在硬件、固件、操作系统、网络、AI 推理、Web 与语音之间的
碎片串成一条可运行、可调试的链路，是这门课程给我最大的工程收获。它
让我对"嵌入式系统"这四个字的理解从"在板子上跑一段代码"扩展到了
"让传感、算力、前端和用户在同一条反馈回路里协同工作"。也是从这门
课开始，我真正具备了去独立承担一个端到端嵌入式 AI 项目的工程能力，
能够在复杂依赖关系下做出取舍、定位问题、并把系统逐步推向稳定可用。

回过头看，IronBuddy 还不是一个可直接产品化的成品，但它已经回答了
最初提出的问题：当视觉与肌电互相补位时，居家健身系统可以从看动作
姿态走向理解动作背后的发力质量；而这条路要走通，需要的远不止一个
模型，更是一整套硬件、嵌入式软件、AI 推理与人机交互协同设计的能力。
